\section{Discussione}

I risultati indicano che i blocchi ricevuti sono il segnale comunitario piu' informativo tra quelli analizzati, ma il loro potere dipende dall'esposizione pubblica dell'account. La distinzione centrale e' tra due regimi.

Nel regime degli account esposti, cioe' 0p, p0 e soprattutto pp, gli account sono osservabili dalla comunita', possono accumulare blocchi e mostrano un segnale temporalmente concentrato vicino al takedown. La Random Forest non produce una previsione completa, ma distingue parzialmente positivi e negativi: la precision e' alta, mentre la recall resta moderata. Questo profilo e' coerente con un segnale utile quando presente, ma incapace di coprire l'intera popolazione dei takedown. Il gruppo pp, che combina attivita' e visibilita', e' quello con maggiore separabilita' predittiva.

Nel regime del corner 00, molti positivi hanno zero post e zero follow ricevuti nella finestra osservata. Essi accumulano pochi blocchi in valore assoluto e spesso presentano lifetime molto breve. Questa dinamica e' coerente con account usa-e-getta, spam, ban evasion, impersonation o altri pattern anti-abuso account-level, ma il dataset non osserva direttamente la causa del takedown e non consente una classificazione sostantiva dei singoli account. L'interpretazione prudente e' che una quota rilevante di takedown ufficiali avvenga prima che un segnale comunitario pubblico possa formarsi.

La temporalita' del segnale e' particolarmente rilevante. Nei gruppi esposti, il giorno immediatamente precedente all'evento ha un peso elevato sia descrittivamente sia nella feature importance. Questo suggerisce una capacita' predittiva ravvicinata: i blocchi non anticipano necessariamente il takedown su orizzonti lunghi, ma contengono informazione nelle ore e nei giorni immediatamente precedenti. La previsione e' quindi utile in senso statistico e operativo limitato, non come sistema generale di early warning distante nel tempo.

Labels e modlist mostrano un risultato diverso. Alcuni labeler terzi hanno lift elevati, ma la coverage e' troppo bassa per sostenere una predizione sistematica. Questo non implica che i labeler siano teoricamente irrilevanti: possono servire nicchie specifiche, policy particolari o comunita' delimitate. Nel dataset operativo, pero', non coprono una quota sufficiente dei positivi pre-takedown. La stessa conclusione vale per le modlist, che rappresentano un segnale comunitario aggregato interessante ma raro nella finestra considerata.

Nel complesso, la moderazione comunitaria e quella ufficiale non appaiono come sistemi sovrapposti perfettamente. Per gli account esposti, i blocchi possono precedere temporalmente l'enforcement ufficiale e contenere informazione predittiva. Per molti account 00, l'enforcement ufficiale sembra invece agire su segnali account-level o anti-abuso non osservati pubblicamente. Il contributo principale del lavoro e' rendere empiricamente visibile questa eterogeneita'.
